\documentclass{article}
\usepackage[T1]{fontenc}
\usepackage{lmodern}
\usepackage{graphicx}
\usepackage{amsmath,amssymb}
\usepackage[a4paper,margin=2cm]{geometry}
\usepackage[absolute,overlay]{textpos}
\setlength{\TPHorizModule}{1mm}
\setlength{\TPVertModule}{1mm}
\usepackage[indonesian]{babel}
\usepackage{hyperref}
\setlength{\emergencystretch}{2em}
% unit: Halaman 11

\begin{document}

% === Halaman 11 (python/native) ===

11

Contoh pembangkitan kunci oleh Alice

• Misalkan Alice ingin menghitung pasangan kunci publik dan kunci privatnya. 

• Misalkan Alice memilih p = 47 dan q = 71 (keduanya prima), maka dapat 

dihitung: 


n = p  q = 3337 


(n) = (p – 1)(q – 1) = (47 – 1)(71 – 1) = 46 x 60 = 3220. 

• Alice memilih kunci publik e = 79 (yang relatif prima dengan 3220 karena 

pembagi bersama terbesarnya adalah 1).  

• Kunci publik Alice adalah (e, n) = (79, 3337). Nilai e dan n dapat dipublikasikan 

kepada umum.

\end{document}
